\chapter{Epidural Steroid Injection}

\section*{Overview}
An epidural steroid injection delivers corticosteroid into the epidural space around the spinal nerve roots to treat radicular pain from disc herniation, stenosis, or inflammation.

\section*{Alternative Names}
Epidural injection; transforaminal epidural block; interlaminar epidural steroid injection; ESI

\section*{Principle}
Corticosteroid is injected into the epidural space, where it reduces inflammation and swelling around irritated nerve roots, relieving radicular pain. Local anesthetic may be added for immediate relief and diagnosis.

\section*{History}
Epidural steroid injections for sciatica were introduced in the 1950s and became one of the most commonly performed interventional pain procedures, refined with fluoroscopic guidance in recent decades.

\section*{Why It Is Done}
Performed for radicular back or neck pain with nerve root irritation from disc herniation, spinal stenosis, or postoperative scarring that has not responded to conservative therapy.

\section*{Is This a Primary or Secondary Test or Procedure}
Secondary treatment for radicular pain; primary interventional therapy when conservative care fails.

\section*{How It Is Performed}
Performed in an outpatient setting with fluoroscopic guidance. The patient lies prone, the level is identified, and a needle is advanced into the epidural space using a transforaminal, interlaminar, or caudal approach. Contrast injection confirms placement, and corticosteroid is injected.

\section*{Preparation}
Anticoagulants are held per guidelines, informed consent is obtained, and patients are observed briefly after the procedure. Fasting is not typically required.

\section*{Contraindications and Precautions}
Active infection, bleeding disorder, anticoagulation, allergy to injectate, and local malignancy are contraindications. Precaution: transforaminal cervical and high lumbar injections carry a risk of spinal cord injury.

\section*{Risks}
Transient pain flare, bleeding, infection, dural puncture with headache, nerve injury, increased blood glucose, and very rarely spinal cord injury or infarction.

\section*{Aftercare and Recovery}
Patients may experience immediate relief from local anesthetic and delayed steroid effect over 3-7 days. Activity is limited briefly, and a series of injections may be planned.

\section*{Outcome and Follow-up}
Relief of radicular pain for weeks to months confirms an inflammatory nerve root component; the response guides repeat treatment or surgical referral.

\section*{Expected Results}
Significant reduction in radicular pain and improved function lasting weeks to months without complication.

\section*{Follow-up}
Clinical reassessment after each injection; imaging or surgical evaluation if no durable benefit is obtained after a limited series.

\section*{References}
\begin{enumerate}
  \item MedlinePlus. Epidural steroid injection. U.S. National Library of Medicine; 2025.
  \item Current Medical Diagnosis and Treatment. McGraw-Hill; 2025.
\end{enumerate}

\clearpage
